GPT-4o miniや小型系列の更新は、最大モデルの性能競争とは別に、cost、latency、throughput、localityを含む『能力密度』の競争を前景化した。小型化や長文脈対応は、それ自体で大型モデルと同等の品質を意味しないが、用途ごとのmodel choiceを増やし、API/edge/localの設計余地を広げた。\autocite{src-031c86eaad2b,src-ec05e05a8071,src-f086061a9e86}


FlashAttention-3のようなkernel/runtime側の進歩は、model architectureそのものを変えずにmemory trafficや実行効率へ影響する。一方でcaching、distillation、efficient generationに関する研究は、モデルの能力とserving costの間に複数の最適化レイヤがあることを示す。model releaseの性能とruntime optimizationの効果を同じ数字へ畳み込まないことが重要である。\autocite{src-3aa7f9dab810,src-ecac5d71389f,src-ee43656ed6da,src-fc4f100c5e5a}


LTX-Videoのようなefficient media generationも同じ延長線上にある。video生成ではqualityだけでなく、sampling cost、VRAM、resolution、duration、local execution feasibilityが実用性を大きく左右するため、生成結果の見栄えだけでなくdeployment条件を技術仕様として扱う必要がある。\autocite{src-fd5d0edaf394}


